% GENERATED FILE. Edit canonical lessons, then run npm run generate:books.
% canonical-source-hash: fnv1a64:5c8c4c9992112ced
% canonical-lessons: PT-C478-nostalgia, PT-C478-melancolia, PT-C478-solidao, PT-C478-solitario, PT-C478-desgosto

\chapter{Solidão e saudade}
\label{ch:pt-solidao-e-saudade}
\hlchaptermodality{478}

\begin{chapteropening}
\textbf{By the end of this chapter:} \emph{I can talk about nostalgia, loneliness and sorrow.}

Complete the last lesson of chapter 478: I can talk about nostalgia, loneliness and sorrow.
\end{chapteropening}

\section[a nostalgia]{a nostalgia — nostalgia}

\label{lesson:PT-C478-nostalgia}



\begin{quote}
Before the new one: say the Portuguese for sigh, then the Portuguese for burst of laughter.
\end{quote}

\subsection*{You'll want to know: a nostalgia}
\textbf{a nostalgia} — \textquotedblleft{}nostalgia\textquotedblright{}.

\begin{quote}
A avó olha para as fotografias antigas com nostalgia.
\end{quote}

\begin{quote}
\emph{Grandma looks at the old photographs with nostalgia.}
\end{quote}

Close to \emph{a saudade}, but more about a time that is gone.

\begin{cousinweb}
\emph{nostalgia} comes from New Latin \emph{nostalgia} (EtymDB), a modern coinage from Greek \emph{nóstos} \textquotedblleft{}return home\textquotedblright{} and \emph{álgos} \textquotedblleft{}pain\textquotedblright{}.
\end{cousinweb}

\begin{grammarlens}[title={What you've built}]
One more word for telling, explaining and describing at B1.
\end{grammarlens}

\subsection*{Your turn}
\begin{itemize}
\raggedright
  \item \emph{Say it:} \emph{a nostalgia}
  \item \emph{Say it:} \emph{a nostalgia}, once more
  \item \emph{Say it:} a sentence of your own with \emph{a nostalgia}
  \item \emph{Recall:} say the Portuguese for sigh, then the Portuguese for burst of laughter, then say \emph{a nostalgia} again
\end{itemize}

\begin{tcolorbox}[breakable,colback=teal!4,colframe=teal!35!black,title={Before you move on}]
What does a nostalgia mean? (\textquotedblleft{}Nostalgia\textquotedblright{}.) Say it once more.
\end{tcolorbox}
\section[a melancolia]{a melancolia — melancholy}

\label{lesson:PT-C478-melancolia}



\begin{quote}
Before the new one: say the Portuguese for burst of laughter, then the Portuguese for nostalgia.
\end{quote}

\subsection*{You'll want to know: a melancolia}
\textbf{a melancolia} — \textquotedblleft{}melancholy\textquotedblright{}.

\begin{quote}
Os dias de chuva enchem-me de melancolia.
\end{quote}

\begin{quote}
\emph{Rainy days fill me with melancholy.}
\end{quote}

The adjective is \emph{melancólico}.

\begin{cousinweb}
\emph{melancolia} goes back, through Old Portuguese \emph{melanconia} (EtymDB), to Latin \emph{melancholia} (Lewis \& Short), from Greek \emph{melankholía} \textquotedblleft{}black bile\textquotedblright{}, which old medicine blamed for sadness.
\end{cousinweb}

\begin{grammarlens}[title={What you've built}]
One more word for telling, explaining and describing at B1.
\end{grammarlens}

\subsection*{Your turn}
\begin{itemize}
\raggedright
  \item \emph{Say it:} \emph{a melancolia}
  \item \emph{Say it:} \emph{a melancolia}, once more
  \item \emph{Say it:} a sentence of your own with \emph{a melancolia}
  \item \emph{Recall:} say the Portuguese for burst of laughter, then the Portuguese for nostalgia, then say \emph{a melancolia} again
\end{itemize}

\begin{tcolorbox}[breakable,colback=teal!4,colframe=teal!35!black,title={Before you move on}]
What does a melancolia mean? (\textquotedblleft{}Melancholy\textquotedblright{}.) Say it once more.
\end{tcolorbox}
\section[a solidão]{a solidão — loneliness, solitude}

\label{lesson:PT-C478-solidao}



\begin{quote}
Before the new one: say the Portuguese for nostalgia, then the Portuguese for melancholy.
\end{quote}

\subsection*{You'll want to know: a solidão}
\textbf{a solidão} — \textquotedblleft{}loneliness, solitude\textquotedblright{}.

\begin{quote}
Muitos idosos vivem em solidão.
\end{quote}

\begin{quote}
\emph{Many elderly people live in loneliness.}
\end{quote}

\emph{Sentir solidão} — to feel lonely.

\begin{cousinweb}
\emph{solidão} is inherited from Latin \emph{sōlitūdō} \textquotedblleft{}loneliness, a lonely place\textquotedblright{} (EtymDB), built on \emph{sōlus} \textquotedblleft{}alone\textquotedblright{}.
\end{cousinweb}

\begin{grammarlens}[title={What you've built}]
One more word for telling, explaining and describing at B1.
\end{grammarlens}

\subsection*{Your turn}
\begin{itemize}
\raggedright
  \item \emph{Say it:} \emph{a solidão}
  \item \emph{Say it:} \emph{a solidão}, once more
  \item \emph{Say it:} a sentence of your own with \emph{a solidão}
  \item \emph{Recall:} say the Portuguese for nostalgia, then the Portuguese for melancholy, then say \emph{a solidão} again
\end{itemize}

\begin{tcolorbox}[breakable,colback=teal!4,colframe=teal!35!black,title={Before you move on}]
What does a solidão mean? (\textquotedblleft{}Loneliness, solitude\textquotedblright{}.) Say it once more.
\end{tcolorbox}
\section[solitário]{solitário — lonely, solitary}

\label{lesson:PT-C478-solitario}



\begin{quote}
Before the new one: say the Portuguese for melancholy, then the Portuguese for loneliness.
\end{quote}

\subsection*{You'll want to know: solitário}
\textbf{solitário} — \textquotedblleft{}lonely, solitary\textquotedblright{}.

\begin{quote}
É um homem solitário que vive no campo.
\end{quote}

\begin{quote}
\emph{He is a solitary man who lives in the country.}
\end{quote}

Also a noun: \emph{um solitário} — a loner.

\begin{cousinweb}
\emph{solitário} is a learned borrowing from Latin \emph{sōlitārius} \textquotedblleft{}lonely, solitary\textquotedblright{} (EtymDB), from \emph{sōlus} \textquotedblleft{}alone\textquotedblright{}.
\end{cousinweb}

\begin{grammarlens}[title={What you've built}]
One more word for telling, explaining and describing at B1.
\end{grammarlens}

\subsection*{Your turn}
\begin{itemize}
\raggedright
  \item \emph{Say it:} \emph{solitário}
  \item \emph{Say it:} \emph{solitário}, once more
  \item \emph{Say it:} a sentence of your own with \emph{solitário}
  \item \emph{Recall:} say the Portuguese for melancholy, then the Portuguese for loneliness, then say \emph{solitário} again
\end{itemize}

\begin{tcolorbox}[breakable,colback=teal!4,colframe=teal!35!black,title={Before you move on}]
What does solitário mean? (\textquotedblleft{}Lonely, solitary\textquotedblright{}.) Say it once more.
\end{tcolorbox}
\section[o desgosto]{o desgosto — grief, sorrow; heartbreak}

\label{lesson:PT-C478-desgosto}



\begin{quote}
Before the new one: say the Portuguese for loneliness, then the Portuguese for lonely.
\end{quote}

\subsection*{You'll want to know: o desgosto}
\textbf{o desgosto} — \textquotedblleft{}grief, sorrow; heartbreak\textquotedblright{}.

\begin{quote}
A morte do cão foi um grande desgosto para as crianças.
\end{quote}

\begin{quote}
\emph{The dog's death was a great sorrow for the children.}
\end{quote}

Stronger than \emph{tristeza}. \emph{Ter um desgosto} — to suffer a heartbreak.

\begin{cousinweb}
\emph{desgosto} is built from \emph{des-} and \emph{gosto} \textquotedblleft{}taste, pleasure\textquotedblright{}, which is inherited from Latin \emph{gustus} \textquotedblleft{}a tasting, taste\textquotedblright{} (Lewis \& Short).
\end{cousinweb}

\begin{grammarlens}[title={What you've built}]
One more word for telling, explaining and describing at B1.
\end{grammarlens}

\subsection*{Your turn}
\begin{itemize}
\raggedright
  \item \emph{Say it:} \emph{o desgosto}
  \item \emph{Say it:} \emph{o desgosto}, once more
  \item \emph{Say it:} a sentence of your own with \emph{o desgosto}
  \item \emph{Recall:} say the Portuguese for loneliness, then the Portuguese for lonely, then say \emph{o desgosto} again
\end{itemize}

\begin{tcolorbox}[breakable,colback=teal!4,colframe=teal!35!black,title={Before you move on}]
What does o desgosto mean? (\textquotedblleft{}Grief, sorrow; heartbreak\textquotedblright{}.) Say it once more.
\end{tcolorbox}
